% GENERATED FILE. Edit canonical lessons, then run npm run generate:books.
% canonical-source-hash: fnv1a64:913c38d13225c347
% canonical-lessons: PA-C311-bhalna, PA-C311-chor, PA-C311-nassna, PA-C311-khohna, PA-C311-nuksan

\chapter{Gone Missing}
\label{ch:pa-a2-gone-missing}
\hlchaptermodality{311}

\begin{chapteropening}
\textbf{By the end of this chapter:} \emph{I can say what went missing and who took it.}

Complete the last lesson of chapter 311: I can say what went missing and who took it.
\end{chapteropening}

\section[\texorpdfstring{\pa{ਭਾਲਣਾ} (bhālṇā)}{ਭਾਲਣਾ (bhālṇā)}]{\pa{ਭਾਲਣਾ} (bhālṇā) — to seek; to hunt for}

\label{lesson:PA-C311-bhalna}



\begin{quote}
Before the new one: say the Punjabi for to break something, then the Punjabi for to get spilt.
\end{quote}

\subsection*{You'll want to know: \pa{ਭਾਲਣਾ}}
\textbf{\pa{ਭਾਲਣਾ}} — \emph{bhālṇā} — \textquotedblleft{}to seek; to hunt for\textquotedblright{}.

\begin{quote}
\pa{ਕੱਲ੍ਹ} \pa{ਮੈਂ} \pa{ਬਾਜ਼ਾਰ} \pa{ਵਿੱਚ} \pa{ਆਪਣੀ} \pa{ਚਾਬੀ} \pa{ਭਾਲੀ।}
\end{quote}

\begin{quote}
\emph{kall maiṁ bāzār vich āpṇī chābī bhālī.}
\end{quote}

\begin{quote}
Yesterday I hunted for my key in the market.
\end{quote}

Past (I …): \pa{ਮੈਂ} \pa{ਕਮਰੇ} \pa{ਵਿੱਚ} \pa{ਆਪਣੀ} \pa{ਕਿਤਾਬ} \pa{ਭਾਲੀ}, \emph{maiṁ kamre vich āpṇī kitāb bhālī}.

\begin{cousinweb}
\textbf{\pa{ਭਾਲਣਾ}} (\emph{bhālṇā}) is inherited from Sanskrit \emph{bhālayate}, \textquotedblleft{}describes, beholds\textquotedblright{}; CDIAL lists Punjabi \emph{bhālṇā}, \textquotedblleft{}to seek\textquotedblright{}, under it (CDIAL 9474).
\end{cousinweb}

\begin{grammarlens}[title={What you've built}]
One more word for saying what happened.
\end{grammarlens}

\subsection*{Your turn}
\begin{itemize}
\raggedright
  \item \emph{Say it:} \emph{bhālṇā}
  \item \emph{Say it:} \emph{bhālṇā}, once more
  \item \emph{Say it:} \emph{kall maiṁ bāzār vich āpṇī chābī bhālī.} — Yesterday I hunted for my key in the market.
  \item \emph{Recall:} say the Punjabi for to break something, then the Punjabi for to get spilt, then say \emph{bhālṇā} again
\end{itemize}

\begin{tcolorbox}[breakable,colback=teal!4,colframe=teal!35!black,title={Before you move on}]
What does \pa{ਭਾਲਣਾ} mean? (\textquotedblleft{}To seek; to hunt for\textquotedblright{}.) Say it once more.
\end{tcolorbox}
\section[\texorpdfstring{\pa{ਚੋਰ} (chor)}{ਚੋਰ (chor)}]{\pa{ਚੋਰ} (chor) — a thief}

\label{lesson:PA-C311-chor}



\begin{quote}
Before the new one: say the Punjabi for to get spilt, then the Punjabi for to seek.
\end{quote}

\subsection*{You'll want to know: \pa{ਚੋਰ}}
\textbf{\pa{ਚੋਰ}} — \emph{chor} — \textquotedblleft{}a thief\textquotedblright{}.

\begin{quote}
\pa{ਕੱਲ੍ਹ} \pa{ਰਾਤ} \pa{ਇੱਕ} \pa{ਚੋਰ} \pa{ਸਾਡੇ} \pa{ਘਰ} \pa{ਆਇਆ।}
\end{quote}

\begin{quote}
\emph{kall rāt ikk chor sāḍe ghar āiā.}
\end{quote}

\begin{quote}
A thief came to our house last night.
\end{quote}

A masculine noun.

\begin{cousinweb}
\textbf{\pa{ਚੋਰ}} (\emph{chor}) is inherited from Sanskrit \emph{cora}, \textquotedblleft{}a thief\textquotedblright{}; CDIAL lists Punjabi \emph{cor}, \textquotedblleft{}thief\textquotedblright{}, under it (CDIAL 4931).
\end{cousinweb}

\begin{grammarlens}[title={What you've built}]
One more word for saying what happened.
\end{grammarlens}

\subsection*{Your turn}
\begin{itemize}
\raggedright
  \item \emph{Say it:} \emph{chor}
  \item \emph{Say it:} \emph{chor}, once more
  \item \emph{Say it:} \emph{kall rāt ikk chor sāḍe ghar āiā.} — A thief came to our house last night.
  \item \emph{Recall:} say the Punjabi for to get spilt, then the Punjabi for to seek, then say \emph{chor} again
  \item \emph{Write it:} copy \textbf{\pa{ਚੋਰ}} once, checking the whole word
\end{itemize}

\begin{tcolorbox}[breakable,colback=teal!4,colframe=teal!35!black,title={Before you move on}]
What does \pa{ਚੋਰ} mean? (\textquotedblleft{}A thief\textquotedblright{}.) Say it once more.
\end{tcolorbox}
\section[\texorpdfstring{\pa{ਨੱਸਣਾ} (nassṇā)}{ਨੱਸਣਾ (nassṇā)}]{\pa{ਨੱਸਣਾ} (nassṇā) — to run away}

\label{lesson:PA-C311-nassna}



\begin{quote}
Before the new one: say the Punjabi for to seek, then the Punjabi for a thief.
\end{quote}

\subsection*{You'll want to know: \pa{ਨੱਸਣਾ}}
\textbf{\pa{ਨੱਸਣਾ}} — \emph{nassṇā} — \textquotedblleft{}to run away\textquotedblright{}.

\begin{quote}
\pa{ਚੋਰ} \pa{ਰਾਤ} \pa{ਨੂੰ} \pa{ਪਿਛਲੇ} \pa{ਦਰਵਾਜ਼ੇ} \pa{ਤੋਂ} \pa{ਨੱਸਿਆ।}
\end{quote}

\begin{quote}
\emph{chor rāt nū̃ pichhle darvāze ton nassiā.}
\end{quote}

\begin{quote}
The thief ran away through the back door at night.
\end{quote}

Past (I …): \pa{ਮੈਂ} \pa{ਕੁੱਤੇ} \pa{ਤੋਂ} \pa{ਨੱਸਿਆ}, \emph{maiṁ kutte ton nassiā}.

\begin{cousinweb}
\textbf{\pa{ਨੱਸਣਾ}} (\emph{nassṇā}) is inherited from Sanskrit \emph{naśyati}, \textquotedblleft{}is lost, perishes, runs away\textquotedblright{}; CDIAL lists Punjabi \emph{nassṇā}, \textquotedblleft{}to flee\textquotedblright{}, under it (CDIAL 7027).
\end{cousinweb}

\begin{grammarlens}[title={What you've built}]
One more word for saying what happened.
\end{grammarlens}

\subsection*{Your turn}
\begin{itemize}
\raggedright
  \item \emph{Say it:} \emph{nassṇā}
  \item \emph{Say it:} \emph{nassṇā}, once more
  \item \emph{Say it:} \emph{chor rāt nū̃ pichhle darvāze ton nassiā.} — The thief ran away through the back door at night.
  \item \emph{Recall:} say the Punjabi for to seek, then the Punjabi for a thief, then say \emph{nassṇā} again
  \item \emph{Write it:} copy \textbf{\pa{ਨੱਸਣਾ}} once, checking the whole word
\end{itemize}

\begin{tcolorbox}[breakable,colback=teal!4,colframe=teal!35!black,title={Before you move on}]
What does \pa{ਨੱਸਣਾ} mean? (\textquotedblleft{}To run away\textquotedblright{}.) Say it once more.
\end{tcolorbox}
\section[\texorpdfstring{\pa{ਖੋਹਣਾ} (khohṇā)}{ਖੋਹਣਾ (khohṇā)}]{\pa{ਖੋਹਣਾ} (khohṇā) — to snatch}

\label{lesson:PA-C311-khohna}



\begin{quote}
Before the new one: say the Punjabi for a thief, then the Punjabi for to run away.
\end{quote}

\subsection*{You'll want to know: \pa{ਖੋਹਣਾ}}
\textbf{\pa{ਖੋਹਣਾ}} — \emph{khohṇā} — \textquotedblleft{}to snatch\textquotedblright{}.

\begin{quote}
\pa{ਮੈਂ} \pa{ਬੱਚੇ} \pa{ਦੇ} \pa{ਹੱਥ} \pa{ਤੋਂ} \pa{ਚਾਕੂ} \pa{ਖੋਹਿਆ।}
\end{quote}

\begin{quote}
\emph{maiṁ bachche de hatth ton chākū khohiā.}
\end{quote}

\begin{quote}
I snatched the knife out of the child's hand.
\end{quote}

Past (I …): \pa{ਮੈਂ} \pa{ਕੁੱਤੇ} \pa{ਤੋਂ} \pa{ਰੋਟੀ} \pa{ਖੋਹੀ}, \emph{maiṁ kutte ton roṭī khohī}.

\begin{cousinweb}
\textbf{\pa{ਖੋਹਣਾ}} (\emph{khohṇā}) is inherited from a reconstructed \emph{*skoṣati}, \textquotedblleft{}plucks out, pokes\textquotedblright{}; CDIAL lists Punjabi \emph{khohṇā}, \textquotedblleft{}to take away by force\textquotedblright{}, under it (CDIAL 13661).
\end{cousinweb}

\begin{grammarlens}[title={What you've built}]
One more word for saying what happened.
\end{grammarlens}

\subsection*{Your turn}
\begin{itemize}
\raggedright
  \item \emph{Say it:} \emph{khohṇā}
  \item \emph{Say it:} \emph{khohṇā}, once more
  \item \emph{Say it:} \emph{maiṁ bachche de hatth ton chākū khohiā.} — I snatched the knife out of the child's hand.
  \item \emph{Recall:} say the Punjabi for a thief, then the Punjabi for to run away, then say \emph{khohṇā} again
  \item \emph{Write it:} copy \textbf{\pa{ਖੋਹਣਾ}} once, checking the whole word
\end{itemize}

\begin{tcolorbox}[breakable,colback=teal!4,colframe=teal!35!black,title={Before you move on}]
What does \pa{ਖੋਹਣਾ} mean? (\textquotedblleft{}To snatch\textquotedblright{}.) Say it once more.
\end{tcolorbox}
\section[\texorpdfstring{\pa{ਨੁਕਸਾਨ} (nuksān)}{ਨੁਕਸਾਨ (nuksān)}]{\pa{ਨੁਕਸਾਨ} (nuksān) — damage; a loss}

\label{lesson:PA-C311-nuksan}



\begin{quote}
Before the new one: say the Punjabi for to run away, then the Punjabi for to snatch.
\end{quote}

\subsection*{You'll want to know: \pa{ਨੁਕਸਾਨ}}
\textbf{\pa{ਨੁਕਸਾਨ}} — \emph{nuksān} — \textquotedblleft{}damage; a loss\textquotedblright{}.

\begin{quote}
\pa{ਉਸ} \pa{ਹਾਦਸੇ} \pa{ਵਿੱਚ} \pa{ਬਹੁਤ} \pa{ਨੁਕਸਾਨ} \pa{ਹੋਇਆ।}
\end{quote}

\begin{quote}
\emph{us hādse vich bahut nuksān hoiā.}
\end{quote}

\begin{quote}
There was a lot of damage in that accident.
\end{quote}

A masculine noun.

\begin{cousinweb}
\textbf{\pa{ਨੁਕਸਾਨ}} (\emph{nuksān}) is a loan, usually traced through Persian to Arabic \emph{nuqṣān}, \textquotedblleft{}a decrease, a shortfall\textquotedblright{}.
\end{cousinweb}

\begin{grammarlens}[title={What you've built}]
One more word for saying what happened.
\end{grammarlens}

\subsection*{Your turn}
\begin{itemize}
\raggedright
  \item \emph{Say it:} \emph{nuksān}
  \item \emph{Say it:} \emph{nuksān}, once more
  \item \emph{Say it:} \emph{us hādse vich bahut nuksān hoiā.} — There was a lot of damage in that accident.
  \item \emph{Recall:} say the Punjabi for to run away, then the Punjabi for to snatch, then say \emph{nuksān} again
  \item \emph{Write it:} copy \textbf{\pa{ਨੁਕਸਾਨ}} once, checking the whole word
\end{itemize}

\begin{tcolorbox}[breakable,colback=teal!4,colframe=teal!35!black,title={Before you move on}]
What does \pa{ਨੁਕਸਾਨ} mean? (\textquotedblleft{}Damage; a loss\textquotedblright{}.) Say it once more.
\end{tcolorbox}
